\chapter{Conclusion}
\label{ch:conclusion}

This thesis built continuous $1$~km gridded daily precipitation
fields over northern and central Germany from the ReKIS station
network. Three families of estimator were implemented on a common
preprocessing pipeline and evaluated on the same held-out subset of
the network: Ordinary Kriging as the geostatistical baseline,
gradient boosting as a flexible regressor over engineered covariates,
and Bayesian Neural Fields as a probabilistic spatiotemporal model.

The three objectives stated in Chapter~\ref{ch:introduction} were
met as follows. The first objective, implementation and validation
of Ordinary Kriging as a baseline, was fulfilled by the two-stage
indicator-plus-amount pipeline of Chapter~\ref{ch:kriging}, whose
leave-one-out cross-validation under MAE and CRPS established the
geostatistical reference against which the remaining methods are
measured. The second objective, assessment of terrain elevation
as an auxiliary predictor, was addressed in the same chapter
through the elevation-aware monthly normals built on the Copernicus
DEM and confirmed that orographic information reduces the residual
trend prior to spatial interpolation. The third objective,
contextualisation against machine-learning alternatives, was met
by the gradient boosting pipeline of Chapter~\ref{ch:gbm}
and the Bayesian Neural Field of Chapter~\ref{ch:bayesnf}, both
trained and validated on the same data splits as the geostatistical
baseline; the cross-method comparison of Chapter~\ref{ch:comparison}
quantified the trade-off between interpretability and predictive
accuracy and provided the empirical basis for selecting the
Bayesian Neural Field as the production model.

The cross-method evaluation of Chapter~\ref{ch:comparison} confirmed
that the Bayesian Neural Field attains the lowest holdout MAE and CRPS
among the three methods and is the only one of them that returns a
calibrated predictive distribution at every grid cell; gradient
boosting placed between Ordinary Kriging and the Bayesian Neural Field
on both metrics. The selected model was then materialised as a
four-resolution gridded daily precipitation product for the worked
example month and distributed as a Parquet dataset over the four
nested target resolutions.

Inspection of the resulting fields confirmed that the high-resolution
Copernicus DEM and the SoilGrids covariates are a necessary input for
resolving local rainfall: without elevation and soil information at
$1$~km, the orographic and surface-controlled components of local
rainfall collapse into the broad station-distance trend, and the
progressive smoothing observed at coarser resolutions is the direct
expression of this mechanism.

Two limitations of the present work point to natural continuations.
The \texttt{bayesnf} library exposes its model and its variational
inference loop at a fairly low level, so that using the framework to
its full strength requires fluent JAX and a working knowledge of
variational inference and Bayesian model construction; at the present
stage of the author's expertise this is the binding constraint on what
Chapter~\ref{ch:bayesnf} could explore, and the hyperparameter sweep
reported there should be read as a first pass over the configuration
space rather than as a settled tuning. A more principled second round
of model tuning, a deeper analysis of the covariate stack already
introduced in Chapter~\ref{ch:bayesnf}, and a direct cross-check of
the gridded product against established national gridded precipitation
datasets would situate the present results among existing operational
references and identify in which regimes and regions the proposed
pipeline improves on or falls behind them.
